\section{Proof-Batching Tools}\label{sec:tools}

% This section presents the two batching tools used repeatedly in the protocol.
% The first handles heterogeneous linear identities whose natural domains
% differ. The second merges many same-form local nonlinear checks into one
% larger virtual instance. Together they cover the two batching patterns that
% recur throughout the rest of the construction.
% Standard divisibility-based zero checks and same-domain batching already appear
% in prior proof systems such as Aurora and CQ++ \cite{aurora19,cq++24}. Our
% contribution here is to isolate two reusable tools: domain-lifting batching for
% heterogeneous linear claims, and interleaving batching for aligned same-form
% nonlinear claims.

\subsection{Domain-Lifting Batching}\label{app:domain-lifting-batch}

As explained earlier, most linear constraints in our construction are first
expressed as zero-checks over some evaluation domain $\mathbb{H}_n$. The
following theorem shows how to batch many such claims together even when their
natural domains differ. The key idea is to lift each claim from its original
domain to a common larger domain, so the lifted claim is still a zero-check
but now over the same domain as the others. After lifting, we can apply an
ordinary same-domain batch to combine all claims into one.

\begin{theorem}[Domain-lifting batching]
\label{thm:domain-lifting-batching}
Let
\[
2 \leq n_1 \leq n_2 \leq \cdots \leq n_m = N
\]
be powers of two. For each $i \in [m]$, let $F_i(X)$ be a virtual polynomial
and suppose the target claim is
\begin{align}
\deg r_i(X) &< n_i-2,\\
F_i(X) &= r_i(X)\cdot X + q_i(X)\nu_{n_i}(X),
\end{align}
for some polynomial $q_i(X)$ and $r_i(X)$.
Because $n_i \mid N$, the ratio $\nu_N(X)/\nu_{n_i}(X)$ is again a polynomial.
Define the lifted polynomial
\[
\widehat{F}_i(X):=
F_i(X)\frac{\nu_N(X)}{\nu_{n_i}(X)}.
\]
Then for a verifier challenge $\eta$, all $m$ claims can be combined into
  the single identity
  \begin{align}\label{eq:domain-lifting-batching}
    \deg r(X) &< N-2,\\
    \sum_{i=1}^{m}\eta^{i-1}\widehat{F}_i(X) &= r(X)\cdot X + q(X)\nu_N(X).
  \end{align}
  And if at least one original claim is false, then
  \eqref{eq:domain-lifting-batching} holds for at most $m-1$ values of $\eta$, i.e. the soundness error is at most $(m-1)/|\mathbb{F}|$.
\end{theorem}

\begin{proof}
Consider there exists $i^*\in [m]$ such that the original
  claim over $F_{i^*}(X)$ is false. 
  This means that $F_{i^*}(X)$ has non-zero
  constant term $c$ when modulo $\nu_{n_{i^*}}(X)$, i.e.
  \begin{align}
    F_{i^*}(X) = r_{i^*}(X)\cdot X + q_{i^*}(X)\nu_{n_{i^*}}(X) + c.
  \end{align}
  For $\widehat{F}_{i^*}(X)$, the constant term will not vanish
  when modulo $\nu_N(X)$, because
  \begin{align}
    \widehat{F}_{i^*}(X) &= F_{i^*}(X)\frac{\nu_N(X)}{\nu_{n_{i^*}}(X)}\\
    &= r_{i^*}(X)\cdot X\frac{\nu_N(X)}{\nu_{n_{i^*}}(X)} + q_{i^*}(X)\nu_N(X) + c\frac{\nu_N(X)}{\nu_{n_{i^*}}(X)}.
  \end{align}
  Since $\nu_N(X)/\nu_{n_{i^*}}(X)$ is a non-zero constant when modulo $\nu_N(X)$, as the last term is the only term contributing to the constant term,
  the constant term of $\widehat{F}_{i^*}(X)$ is also non-zero when modulo $\nu_N(X)$.

Thus, after the random linear combination,
  the constant term remains non-zero except for a probability at most $(m-1)/|\mathbb{F}|$.
\end{proof}

We note two use cases for this batching.

\paragraph{Zero-check.} When $r(X) = 0$ and $r_i(X) = 0$ is required for every $i\in [m]$, the original
claim becomes $F_i(X) = q_i(X)\nu_{n_i}(X)$, which means
$F_i(X)$ vanishes everywhere over $\mathbb{H}_{n_i}$.

\paragraph{Sum-check.} When $r_i(X)$ is a witness committed by the prover,
the original claim becomes $F_i(X) - s_i/n_i = r_i(X)\cdot X + q_i(X)\nu_{n_i}(X)$, which, by the merit of Aurora~\cite{aurora19}, means $F_i(X)$ sums to $s_i$ over $\mathbb{H}_{n_i}$.

We note that we cannot batch zero-check with sum-check due to soundness
considerations. The above theorem only applies to claims of the same form, so we can batch zero-checks together and sum-checks together, but not mix them.


\subsection{Interleaving}\label{app:interleaving}

The intended setting for interleaving batching is that many checks share the
same local rule and the same public parameters, but arrive as separate
PCS commitments and witnesses on possibly different power-of-two domains.
We first record the alignment lemma
used to place smaller objects on larger domains.

\begin{lemma}[Row-replication]
\label{lem:row-replication}
Let $\mathbf{v}\in\mathbb{F}^m$ and let $v(X)$ be its encoding over
$\mathbb{H}_m$. Let $\mathbf{V}\in\mathbb{F}^{r\times m}$ be the matrix whose
every row equals $\mathbf{v}$. Then the encoding polynomial of $\mathbf{V}$ over
$\mathbb{H}_{rm}$ is
\[
V(X)=v(X^r).
\]
\end{lemma}

\begin{proof}
For every $i\in[r]$ and $j\in[m]$,
\[
V(\omega_{rm}^{im+j})
= \mathbf{V}[i,j]
= \mathbf{v}[j]
= v(\omega_m^j)
= v((\omega_{rm}^{im+j})^r),
\]
where the last equality uses $\omega_{rm}^r=\omega_m$. Since both sides have
degree less than $rm$ and agree on $\mathbb{H}_{rm}$, they are identical.
\end{proof}

\begin{definition}[Pairwise interleaving]
\label{def:pairwise-interleaving}
Let $f(X)$ and $g(X)$ be the encoding polynomials of two length-$n$ vectors
over $\mathbb{H}_n$, where $n$ is a power of two. Define
\[
\operatorname{even}_n(X):=\frac{X^n+1}{2},
\qquad
\operatorname{odd}_n(X):=\frac{1-X^n}{2}.
\]
Their evaluations on $\mathbb{H}_{2n}$ are the alternating selectors
$(1,0,1,0,\ldots)$ and $(0,1,0,1,\ldots)$. The pairwise interleaving of
$f(X)$ and $g(X)$ is defined as the virtual polynomial
\[
\operatorname{Int}_n(f,g)(X)
:=
\operatorname{even}_n(X)\,f(X)
+
\operatorname{odd}_n(X)\,g(\omega_{2n}^{-1}X).
\]
Its evaluation vector over $\mathbb{H}_{2n}$ is
\[
(f[0],g[0],f[1],g[1],\ldots,f[n-1],g[n-1]).
\]
\end{definition}

\begin{lemma}[Interleaving]
\label{lem:interleaving}
For each $i\in[m]$, let
$u_i(X)$ be a polynomial encoding $\mathbf{u}_i$
over $\mathbb{H}_{n_i}$, where each
$n_i$ is a power of two. Let $S=\sum_{i\in [m]} n_i$ and
$\log\widehat{n}=\lceil \log_2 S\rceil$. There is a deterministic algorithm that
interleaves these objects into a virtual polynomial $\widehat{u}(X)$ over
$\mathbb{H}_{\widehat{n}}$ and guarantees:
\begin{enumerate}
  \item the evaluations before and after interleaving have the same underlying
  set of entries:
  \[
  \{\widehat{u}(a):a\in\mathbb{H}_{\widehat{n}}\}
  =
  \bigcup_{i=1}^m \{u_i(a):a\in\mathbb{H}_{n_i}\};
  \]
  \item the algorithm runs in $O(m\log \widehat{n})$ time;
  \item any opening claim for $\widehat{u}(X)$ is reduced to
  $m$ openings over $\{u_i\}_{i\in [m]}$, one per polynomial,
  with $O(m\log \widehat{n})$ field operations on the verifier side.
\end{enumerate}
We denote by such algorithm as 
\[ (\widehat{\mathbf{u}}, \widehat{n}) \gets \mathsf{Interleave}((\mathbf{u}_i, n_i)_{i\in [m]}).\]
For simplicity, we may omit the domain sizes when they are clear from context.
\end{lemma}

\begin{proof}
Start from the active multiset $\{(u_i(X),n_i)\}_{i\in[m]}$.
We first store all pairs with same domain size in the same bucket.
Then we perform two scans over the buckets.
In the first scan, process the buckets from small to large domain size and
merge equal-sized pairs whenever a bucket contains at least two active pairs.
Each such merge removes two equal-sized pairs,
uses pairwise interleaving as in
\Cref{def:pairwise-interleaving} to create a new pair,
and inserts the new pair into the appropriate bucket. After this
pass, each bucket contains at most one active pair.

In the second scan, repeatedly take the two active pairs with smallest domain size $n_i \leq n_j$.
If $n_i\neq n_j$, replicate the smaller one by mapping $X\mapsto X^{n_j/n_i}$
as stated in \Cref{lem:row-replication},
then interleave the pair.
Repeat until one active pair $(\widehat{u}(X),\widehat{n})$ remains.

Because the domain sizes are powers of two, the bucket structure contains at most
$\log \widehat{n}$ levels. Each merge removes one active item, and each lift
or interleaving step changes only one level, so the algorithm runs
in $O(m\log \widehat{n})$ time.

For set preservation, lifting is exactly the setting of
Lemma~\ref{lem:row-replication}:
mapping $X\mapsto X^{n_j/n_i}$ repeats
$\mathbf{u}_i$ on the larger domain $\mathbb{H}_{n_j}$ and introduces no new
elements. Clearly, interleaving also introduces no new elements.
As both operations preserve the elements in the vector, repeating them
gives the claimed set equality.

For the final domain size, note that the first scan preserves the total size $S$ and
leaves at most one active pair in each bucket. 
Denote the remaining pairs after the first scan by $(n_i', u_i'(X))$ with
\[
n_1'<\cdots<n_k',
\]
then
\[
\sum_{i=1}^k n_i' = S
\qquad\text{and}\qquad
\log n_k' = \lfloor \log_2 S\rfloor.
\]
If $k=1$, then $S$ is already a power of two and $\widehat{n}=S$. Otherwise
the second scan merges the remaining powers upward and can push the final size
no higher than $2n_k'$. Since $S>n_k'$, we have
\[
\log \widehat{n}= 1 + \log n_k' = \lceil \log_2 S\rceil.
\]
This is exactly the least power of two not smaller than $S$.

For opening reduction, the even/odd selectors are public and sparse, so the verifier can
evaluate them directly. For row replication $u(X)\mapsto u(X^r)$, an opening at
$x$ reduces to an opening of the base polynomial at $x^r$. Hence an opening
claim for the interleaved polynomial reduces to one opening per original
polynomial. The verifier reconstructs the virtual value by traversing the merge
tree and evaluating public selectors, which costs $O(m\log \widehat{n})$ field
operations.
\end{proof}

% \begin{theorem}[Interleaving batching]
% \label{thm:interleaving-batch}
% For each $i\in[m]$, let
% $u_i(X)$ be a polynomial encoding $\mathbf{u}_i$
% over $\mathbb{H}_{n_i}$, where each
% $n_i$ is a power of two. For some local relation $\varphi$,
% the target claim is that
% \[
% \forall i \in [m],\ \forall a\in\mathbb{H}_{n_i}.\ \varphi(u_i(a))=0.
% \]
% There exists an algorithm that combines these $m$ claims into one
% check over one virtual polynomial $\widehat{u}(X)$ and
% $\mathbb{H}_{\widehat{n}}$:
% \[
% \forall a\in \mathbb{H}_{\widehat{n}}.\ \varphi(\widehat{u}(a))=0.
% \]
% The algorithm runs in $O(m\log \widehat{n})$ time
% and guarantees:
% \begin{enumerate}
%   \item all original claims over $u_i$ and $\mathbb{H}_{n_i}$ hold if and only if
%     the claim over $\widehat{u}$ and $\mathbb{H}_{\widehat{n}}$ holds;
%   \item the final domain size satisfies
%   \[
%   \log \widehat{n}
%   =
%   \left\lceil \log_2(n_1+\cdots+n_m)\right\rceil;
%   \]
%   \item any opening claim for $\widehat{u}(X)$ is reduced to
%   $m$ openings over $\{u_i\}_{i\in [m]}$, one per polynomial,
%   with $O(m\log \widehat{n})$ field operations on the verifier side.
% \end{enumerate}
% \end{theorem}

% \begin{proof}
% Apply Lemma~\ref{lem:interleaving} to obtain $\widehat{u}(X)$ over
% $\mathbb{H}_{\widehat{n}}$, where
% $\log\widehat{n}=\lceil\log_2(n_1+\cdots+n_m)\rceil$.
% By the set equality in Lemma~\ref{lem:interleaving}, every value checked in the
% batched instance is one of the original values, and every original value appears
% in the batched instance. Therefore $\varphi$ vanishes on all original domains if
% and only if it vanishes on the interleaved domain. The running-time and opening
% claims are exactly the efficiency guarantees of Lemma~\ref{lem:interleaving}.
% \end{proof}

\subsection{Interleaving Batching}

In our proof system, we use interleaving batching mainly to combine
  two types of check:
  lookup argument and fixed-point quantization check.

\paragraph{Batching lookup.}
We first note that, like~\cite{cq++24}, all table lookups in our construction are
  first reduced to vector lookups.
  It is hence without loss of generality to focus on batching vector lookups.

Suppose we have several lookup families indexed by
$i\in[s]$. For each family, there is a public table $\mathbf{t}_i$ and a
collection of vector queries $\mathbf{f}_{i,1},\ldots,\mathbf{f}_{i,m_i}$ such
that every query satisfies $\mathbf{f}_{i,j}\subseteq \mathbf{t}_i$.

We batch them in three steps. First, for each fixed $i$, we batch all
queries $\mathbf{f}_{i, j}\subseteq \mathbf{t}_i$ 
into one query $\mathbf{f}_i\subseteq \mathbf{t}_i$,
by interleaving $\mathbf{f}_{i, j}$ with \Cref{lem:interleaving}.
Second, we shift each table by a random challenge
$\alpha$ and form $\mathbf{t}'_i:=\alpha\cdot i+\mathbf{t}_i$, so that
different tables are separated with high probability.
We then interleave the
shifted tables into one combined table vector $\mathbf{t}$,
which reduces all lookup instances to $\mathbf{f}_i\subseteq \mathbf{t}$.
Third, we again
interleave $\mathbf{f}_i\subseteq \mathbf{t}$
into one global query vector $\mathbf{f}$.

This yields a single lookup claim
\[
\mathbf{f}\subseteq \mathbf{t},
\]
which can be proved with standard lookup argument.
The benefit is that the prover amortizes the lookup proof across all
tables while keeping the query structure explicit.
The soundness is guaranteed by the challenge $\alpha$,
where note that
\begin{align}
  &\forall i<j, k_1, k_2.\ \Pr[\mathbf{t}_i'[k_1] = \mathbf{t}_j'[k_2]] < 1/|\mathbb{F}|\\
  \implies& \Pr[\exists i<j, k_1, k_2.\ \mathbf{t}_i'[k_1] = \mathbf{t}_j'[k_2]] \leq \frac{s^2\cdot \max_i |\mathbf{t}_i|^2}{|\mathbb{F}|}.
\end{align}
The first equation holds since the elements differ
by a non-zero multiple of $\alpha$, 
and the second is derived from the union bound over the first.

\paragraph{Batching quantizations.}
Suppose given $\mathbf{a}_i\in \mathbb{F}^{n_i}$ and $\mathbf{b}_i\in \mathbb{F}^{n_i}$ for $i\in [m]$,
we want to prove that $\mathbf{a}_i$ is a quantization of $\mathbf{b}_i$, i.e.
\begin{align}
  \forall j\in [n_i].\ \mathbf{a}_i[j] = \left\lfloor \frac{\mathbf{b}_i[j]}{2^q} \right\rfloor
\end{align}
for the same public divisor $2^q$.
Such constraints can be batched easily by interleaving
\begin{align}
  (\mathbf{a}, n) \gets \mathsf{Interleave}(\mathbf{a}_0, \dots, \mathbf{a}_{m-1}),\\
  (\mathbf{b}, n) \gets \mathsf{Interleave}(\mathbf{b}_0, \dots, \mathbf{b}_{m-1}),
\end{align}
and require the proof of
\begin{align}
  \forall j\in [n].\ \mathbf{a}[j] = \left\lfloor \frac{\mathbf{b}[j]}{2^q} \right\rfloor.
\end{align}
Note that as interleaving is deterministic,
 $\mathbf{a}$ is aligned with $\mathbf{b}$, 
 and the batching incurs no extra soundness error.
